\begin{longtable}{|p{\textwidth}|}
\caption{Example CDL description of lat-lon dataset}
\label{tab:cdl-lat-lon} \\
\hline \endfirsthead \endhead \endfoot \hline \endlastfoot
\noindent from: \hfill‘SEA\_SURFACE\_HEIGHT\_day\_mean\_2017-12-29\_ECCO\_V4r4\_latlon\_0p50deg.nc’\hfill\null \\
dimensions: \\
\hline
\rowcolor{YellowGreen}
\hangindent=1.0cm
\hangafter1
time = 1\\
\rowcolor{YellowGreen}
\hangindent=1.0cm
\hangafter1
latitude = 360\\
\rowcolor{YellowGreen}
\hangindent=1.0cm
\hangafter1
longitude = 720\\
\rowcolor{YellowGreen}
\hangindent=1.0cm
\hangafter1
nv = 2\\
\hline
coordinates: \\
\hline
\rowcolor{Apricot}
\hangindent=1.0cm
\hangafter1
\hspace{0.5cm}int32 time (time)\\
\rowcolor{Apricot}
\hangindent=1.5cm
\hangafter1
\hspace{0.5cm}\hspace{0.5cm}time:axis = "T"\\
\rowcolor{Apricot}
\hangindent=1.5cm
\hangafter1
\hspace{0.5cm}\hspace{0.5cm}time:bounds = "time\_bnds"\\
\rowcolor{Apricot}
\hangindent=1.5cm
\hangafter1
\hspace{0.5cm}\hspace{0.5cm}time:coverage\_content\_type = "coordinate"\\
\rowcolor{Apricot}
\hangindent=1.5cm
\hangafter1
\hspace{0.5cm}\hspace{0.5cm}time:long\_name = "center time of averaging period"\\
\rowcolor{Apricot}
\hangindent=1.5cm
\hangafter1
\hspace{0.5cm}\hspace{0.5cm}time:standard\_name = "time"\\
\rowcolor{Apricot}
\hangindent=1.5cm
\hangafter1
\hspace{0.5cm}\hspace{0.5cm}time:units = "hours since 1992-01-01T12:00:00"\\
\rowcolor{Apricot}
\hangindent=1.5cm
\hangafter1
\hspace{0.5cm}\hspace{0.5cm}time:calendar = "proleptic\_gregorian"\\
\rowcolor{Apricot}
\hangindent=1.0cm
\hangafter1
\hspace{0.5cm}float32 latitude (latitude)\\
\rowcolor{Apricot}
\hangindent=1.5cm
\hangafter1
\hspace{0.5cm}\hspace{0.5cm}latitude:axis = "Y"\\
\rowcolor{Apricot}
\hangindent=1.5cm
\hangafter1
\hspace{0.5cm}\hspace{0.5cm}latitude:bounds = "latitude\_bnds"\\
\rowcolor{Apricot}
\hangindent=1.5cm
\hangafter1
\hspace{0.5cm}\hspace{0.5cm}latitude:comment = "uniform grid spacing from -89.75 to 89.75 by 0.5"\\
\rowcolor{Apricot}
\hangindent=1.5cm
\hangafter1
\hspace{0.5cm}\hspace{0.5cm}latitude:coverage\_content\_type = "coordinate"\\
\rowcolor{Apricot}
\hangindent=1.5cm
\hangafter1
\hspace{0.5cm}\hspace{0.5cm}latitude:long\_name = "latitude at grid cell center"\\
\rowcolor{Apricot}
\hangindent=1.5cm
\hangafter1
\hspace{0.5cm}\hspace{0.5cm}latitude:standard\_name = "latitude"\\
\rowcolor{Apricot}
\hangindent=1.5cm
\hangafter1
\hspace{0.5cm}\hspace{0.5cm}latitude:units = "degrees\_north"\\
\rowcolor{Apricot}
\hangindent=1.0cm
\hangafter1
\hspace{0.5cm}float32 longitude (longitude)\\
\rowcolor{Apricot}
\hangindent=1.5cm
\hangafter1
\hspace{0.5cm}\hspace{0.5cm}longitude:axis = "X"\\
\rowcolor{Apricot}
\hangindent=1.5cm
\hangafter1
\hspace{0.5cm}\hspace{0.5cm}longitude:bounds = "longitude\_bnds"\\
\rowcolor{Apricot}
\hangindent=1.5cm
\hangafter1
\hspace{0.5cm}\hspace{0.5cm}longitude:comment = "uniform grid spacing from -179.75 to 179.75 by 0.5"\\
\rowcolor{Apricot}
\hangindent=1.5cm
\hangafter1
\hspace{0.5cm}\hspace{0.5cm}longitude:coverage\_content\_type = "coordinate"\\
\rowcolor{Apricot}
\hangindent=1.5cm
\hangafter1
\hspace{0.5cm}\hspace{0.5cm}longitude:long\_name = "longitude at grid cell center"\\
\rowcolor{Apricot}
\hangindent=1.5cm
\hangafter1
\hspace{0.5cm}\hspace{0.5cm}longitude:standard\_name = "longitude"\\
\rowcolor{Apricot}
\hangindent=1.5cm
\hangafter1
\hspace{0.5cm}\hspace{0.5cm}longitude:units = "degrees\_east"\\
\hline
data variables: \\
\hline
\hangindent=1.0cm
\hangafter1
\hspace{0.5cm}float32 SSH (time, latitude, longitude)\\
\hangindent=1.5cm
\hangafter1
\hspace{0.5cm}\hspace{0.5cm}SSH:\_FillValue = "9.969209968386869e+36"\\
\hangindent=1.5cm
\hangafter1
\hspace{0.5cm}\hspace{0.5cm}SSH:coverage\_content\_type = "modelResult"\\
\hangindent=1.5cm
\hangafter1
\hspace{0.5cm}\hspace{0.5cm}SSH:long\_name = "Dynamic sea surface height anomaly"\\
\hangindent=1.5cm
\hangafter1
\hspace{0.5cm}\hspace{0.5cm}SSH:standard\_name = "sea\_surface\_height\_above\_geoid"\\
\hangindent=1.5cm
\hangafter1
\hspace{0.5cm}\hspace{0.5cm}SSH:units = "m"\\
\hangindent=1.5cm
\hangafter1
\hspace{0.5cm}\hspace{0.5cm}SSH:comment = "Dynamic sea surface height anomaly above the geoid, suitable for comparisons with altimetry sea surface height data products that apply the inverse barometer (IB) correction. Note: SSH is calculated by correcting model sea level anomaly ETAN for three effects: a) global mean steric sea level changes related to density changes in the Boussinesq volume-conserving model (Greatbatch correction, see sterGloH), b) the inverted barometer (IB) effect (see SSHIBC) and c) sea level displacement due to sea-ice and snow pressure loading (see sIceLoad). SSH can be compared with the similarly-named SSH variable in previous ECCO products that did not include atmospheric pressure loading (e.g., Version 4 Release 3). Use SSHNOIBC for comparisons with altimetry data products that do NOT apply the IB correction."\\
\hangindent=1.5cm
\hangafter1
\hspace{0.5cm}\hspace{0.5cm}SSH:coordinates = "time"\\
\hangindent=1.5cm
\hangafter1
\hspace{0.5cm}\hspace{0.5cm}SSH:valid\_min = "-2.4861555099487305"\\
\hangindent=1.5cm
\hangafter1
\hspace{0.5cm}\hspace{0.5cm}SSH:valid\_max = "2.2875382900238037"\\
\hangindent=1.0cm
\hangafter1
\hspace{0.5cm}float32 SSHIBC (time, latitude, longitude)\\
\hangindent=1.5cm
\hangafter1
\hspace{0.5cm}\hspace{0.5cm}SSHIBC:\_FillValue = "9.969209968386869e+36"\\
\hangindent=1.5cm
\hangafter1
\hspace{0.5cm}\hspace{0.5cm}SSHIBC:coverage\_content\_type = "modelResult"\\
\hangindent=1.5cm
\hangafter1
\hspace{0.5cm}\hspace{0.5cm}SSHIBC:long\_name = "The inverted barometer (IB) correction to sea surface height due to atmospheric pressure loading"\\
\hangindent=1.5cm
\hangafter1
\hspace{0.5cm}\hspace{0.5cm}SSHIBC:units = "m"\\
\hangindent=1.5cm
\hangafter1
\hspace{0.5cm}\hspace{0.5cm}SSHIBC:comment = "Not an SSH itself, but a correction to model sea level anomaly (ETAN) required to account for the static part of sea surface displacement by atmosphere pressure loading: SSH = SSHNOIBC - SSHIBC. Note: Use SSH for model-data comparisons with altimetry data products that DO apply the IB correction and SSHNOIBC for comparisons with altimetry data products that do NOT apply the IB correction."\\
\hangindent=1.5cm
\hangafter1
\hspace{0.5cm}\hspace{0.5cm}SSHIBC:coordinates = "time"\\
\hangindent=1.5cm
\hangafter1
\hspace{0.5cm}\hspace{0.5cm}SSHIBC:valid\_min = "-0.5228679180145264"\\
\hangindent=1.5cm
\hangafter1
\hspace{0.5cm}\hspace{0.5cm}SSHIBC:valid\_max = "0.8955588340759277"\\
\hangindent=1.0cm
\hangafter1
\hspace{0.5cm}float32 SSHNOIBC (time, latitude, longitude)\\
\hangindent=1.5cm
\hangafter1
\hspace{0.5cm}\hspace{0.5cm}SSHNOIBC:\_FillValue = "9.969209968386869e+36"\\
\hangindent=1.5cm
\hangafter1
\hspace{0.5cm}\hspace{0.5cm}SSHNOIBC:coverage\_content\_type = "modelResult"\\
\hangindent=1.5cm
\hangafter1
\hspace{0.5cm}\hspace{0.5cm}SSHNOIBC:long\_name = "Sea surface height anomaly without the inverted barometer (IB) correction"\\
\hangindent=1.5cm
\hangafter1
\hspace{0.5cm}\hspace{0.5cm}SSHNOIBC:units = "m"\\
\hangindent=1.5cm
\hangafter1
\hspace{0.5cm}\hspace{0.5cm}SSHNOIBC:comment = "Sea surface height anomaly above the geoid without the inverse barometer (IB) correction, suitable for comparisons with altimetry sea surface height data products that do NOT apply the inverse barometer (IB) correction. Note: SSHNOIBC is calculated by correcting model sea level anomaly ETAN for two effects: a) global mean steric sea level changes related to density changes in the Boussinesq volume-conserving model (Greatbatch correction, see sterGloH), b) sea level displacement due to sea-ice and snow pressure loading (see sIceLoad). In ECCO Version 4 Release 4 the model is forced with atmospheric pressure loading. SSHNOIBC does not correct for the static part of the effect of atmosphere pressure loading on sea surface height (the so-called inverse barometer (IB) correction). Use SSH for comparisons with altimetry data products that DO apply the IB correction."\\
\hangindent=1.5cm
\hangafter1
\hspace{0.5cm}\hspace{0.5cm}SSHNOIBC:coordinates = "time"\\
\hangindent=1.5cm
\hangafter1
\hspace{0.5cm}\hspace{0.5cm}SSHNOIBC:valid\_min = "-2.45104718208313"\\
\hangindent=1.5cm
\hangafter1
\hspace{0.5cm}\hspace{0.5cm}SSHNOIBC:valid\_max = "2.2390522956848145"\\
\end{longtable}